\section{The Vandal Siege and the Death of Augustine}

\subsection{The world closing in}

On 26 September 426, old and planning for his church, Augustine had the
presbyter Eraclius designated as his successor before the assembled people,
the notaries recording address and acclamations alike---the scene described
above, preserved as \work{Letter} 213 (Benedict XVI, General Audience, 16
January 2008). Three years later the frontier
failed. ``A great host of savage foes, Vandals and Alans, with some of the
Gothic tribe interspersed, and various other peoples\ldots\ came by ship
from the regions of Spain across the sea and poured into Africa,'' Possidius
writes as an eyewitness of the aftermath, ``raging with cruelty and
barbarity, they completely devastated everything'' (\work{Vita} 28).
Refugee clergy carried the question to Augustine: might a bishop flee?
His answer to Honoratus of Thiabe, preserved because Possidius copied the
whole letter into the \work{Vita}, ruled that when any of the faithful
remain, the ministers of Christ must remain with them, for the sacraments'
sake---flight is permitted only when the shepherd is personally hunted and
the flock provided for (\work{Letter} 228, in \work{Vita} 30). Peace, he
had written a few months earlier to the count Darius, is the nobler work:
``It is a higher glory still to stay war itself with a word, than to slay
men with the sword, and to procure or maintain peace by peace, not by war''
(\work{Letter} 229.2, as quoted in the General Audience of 16 January
2008).

\subsection{The penitential psalms on the wall}

By May or June 430 the Vandals were before Hippo, one of the last fortified
cities, its walls sheltering refugee bishops---Possidius among them. ``And
lo, in the third month of the siege he succumbed to fever and began to
suffer in his last illness'' (\work{Vita} 29). Possidius's account of the
deathbed became one of the most imitated scenes in Christian biography.
Augustine ``in private conversations frequently told us that even after
baptism had been received exemplary Christians and priests ought not depart
from this life without fitting and appropriate repentance. And this he
himself did\ldots\ For he commanded that the shortest penitential Psalms of
David should be copied for him, and during the days of his sickness as he
lay in bed he would look at these sheets as they hung upon the wall and read
them; and he wept freely and constantly.'' Ten days before the end he asked
that no one enter except at the physicians' visits and mealtimes, ``and he
had all that time free for prayer'' (\work{Vita} 31). He died on 28 August
430, in his seventy-sixth year (Benedict XVI, General Audience, 16 January
2008; Possidius: ``he lived seventy-six years, almost forty of which he
spent as a priest or bishop'').

``He made no will, because as a poor man of God he had nothing from which
to make it. He repeatedly ordered that the library of the church and all the
books should be carefully preserved for future generations'' (\work{Vita}
31). The city fell and burned after the siege; the library survived---the
reason, together with the \work{Retractationes} and the \work{Indiculus},
that more of Augustine remains than of any other ancient author. Possidius
closes with the sentence quoted at the head of this study, counting his
forty years of friendship, and with an observation that has governed
Augustine biography ever since: those who read his works ``profit thereby.
But I believe that they were able to derive greater good from him who heard
and saw him as he spoke in person in the church, and especially those who
knew well his manner of life among men'' (\work{Vita} 31). Of burial he says
only that ``in our presence, after
a service was offered to God for the peaceful repose of his body, he was
buried''---at Hippo, therefore, in 430; everything later belongs to the
history of the relics.
